% =============================================================================
\section{Tarefa 2 --- Ensaios e estimação dos parâmetros}
\label{sec:calib}
% =============================================================================

Cada ensaio foi desenhado para que o valor verdadeiro seja conhecido; a
diferença para o valor medido é o erro do sensor.

\subsection{Ruído e desvio do giroscópio (ensaio A)}

Com o sensor parado a velocidade angular é nula, logo a média de cada eixo do
giroscópio é o desvio $\vect b_g$ e a dispersão em torno dela é o ruído. Para
o acelerómetro, o mesmo registo dá $\sigma_a$. A análise de
Allan~\cite{elsheimy2008} separa o ruído branco das variações lentas do
desvio e indica até que duração compensa integrar.
\todo{valores de $\vect b_g$, $\sigma_g$, $\sigma_a$; comparar com a ficha
técnica.}

\paragraph{Filtragem.} \todo{filtro usado (mediana ou passa-baixo), frequência
de corte e porquê; efeito no ruído; cuidado: filtrar atrasa o sinal, por isso
aplica-se o mesmo filtro aos dois sensores ou um filtro sem atraso
(para a frente e para trás), já que tudo é pós-processado.}

\begin{figure}[H]
\centering
\fig{Photos/ruido.png}{0.8\textwidth}
\caption{\todo{ruído dos dois sensores: desvio de Allan, ou sinal em bruto e
filtrado.}}
\label{fig:ruido}
\end{figure}

\subsection{Acelerómetro: desvio e escala (ensaios B e B2)}

Parado, o acelerómetro só sente a gravidade, que é conhecida. Usaram-se dois
métodos.

\paragraph{Seis posições.} Com um eixo alinhado com a vertical, para cima e
depois para baixo, o sensor devia ler $+g$ e $-g$. Das leituras $m^{+}$ e
$m^{-}$:
\begin{equation}
b = \frac{m^{+} + m^{-}}{2}, \qquad s = \frac{m^{+} - m^{-}}{2g}.
\end{equation}

\paragraph{Posições ao acaso.} Em qualquer orientação parada, o módulo da
aceleração corrigida tem de ser $g$. Com $N$ orientações bem distribuídas,
procuram-se os parâmetros que minimizam~\cite{tedaldi2014}
\begin{equation}
\sum_{k=1}^{N}\Bigl(\bigl\lVert S_a^{-1}(\med a_k - \vect b_a)\bigr\rVert - g\Bigr)^2 .
\end{equation}
Não exige alinhar o sensor, mas também não diz como ele está montado: o
módulo é igual para qualquer rotação.

\paragraph{Com o robô.} As mesmas orientações, impostas pelo robô, são
conhecidas. Isso permite verificar o método anterior e estimar a rotação
$R_{ms}$ entre o sensor e a ferramenta, juntamente com a inclinação da mesa.

\begin{figure}[H]
\centering
\fig{Photos/calib_acc.png}{0.75\textwidth}
\caption{\todo{módulo da aceleração em cada posição, antes e depois de
calibrar: os pontos devem passar a cair em $1\,g$.}}
\label{fig:calibacc}
\end{figure}

\subsection{Giroscópio: escala (ensaio C)}

Uma junta do robô roda a velocidade constante conhecida, nos dois sentidos e
a várias velocidades, com o eixo de rotação alinhado sucessivamente com cada
eixo do sensor. A escala é o declive da reta que relaciona a velocidade medida
com a imposta, e o desvio a ordenada na origem.

\begin{figure}[H]
\centering
\fig{Photos/calib_gyro.png}{0.75\textwidth}
\caption{\todo{velocidade angular medida contra a imposta pelo robô, por eixo,
com a reta ajustada.}}
\label{fig:calibgyro}
\end{figure}

\subsection{Efeito da temperatura (ensaio T)}

\todo{desvio do giroscópio e do acelerómetro em função da temperatura do
integrado, do arranque a frio até estabilizar; declive em
\si{\degree\per\second} por \si{\celsius}; comparar com a ficha técnica.}

\subsection{Resultados}

\begin{table}[H]
\centering
\caption{Parâmetros estimados (média e dispersão de 3 repetições).
\todo{preencher.}}
\label{tab:calib}
\small
\begin{tabular}{llccc}
\toprule
\rowcolor{tabcab}
Parâmetro & Método & $x$ & $y$ & $z$ \\
\midrule
$b_a$ [m\emph{g}] & seis posições & & & \\
                  & posições ao acaso & & & \\
                  & robô & & & \\
$s_a$ [--]        & seis posições & & & \\
                  & posições ao acaso & & & \\
                  & robô & & & \\
$b_g$ [\si{\degree\per\second}] & parado & & & \\
$s_g$ [--]        & robô & & & \\
$\sigma_a$ [m\emph{g}], $\sigma_g$ [\si{\degree\per\second}] & parado & & & \\
\bottomrule
\end{tabular}
\end{table}

\todo{comentário aos resultados: os três métodos concordam? qual é a
dispersão entre repetições? quanto mudou o desvio do giroscópio entre sessões?
orientação do sensor no suporte.}
